\documentclass{article}

\usepackage[utf8]{inputenc}
\usepackage[T1]{fontenc}
\usepackage{hyperref}
\usepackage{tabularx}
\usepackage{array}
\usepackage{fancyhdr}
\usepackage{graphicx}
\usepackage[a4paper]{geometry}
\usepackage{multicol}
\usepackage{listings}
\usepackage{tabto}

\title{Projet de BigData \\ Les bases NoSQL}
\author{par Jordan Baudin, Corentin LeGuen et Geoffrey Spaur}
\date{11 janvier 2018}
\pagestyle{fancy}
\lhead{Projet de BigData - Les bases NoSQL \\ \textbf{M2GIL} - Jordan Baudin, Corentin LeGuen et Geoffrey Spaur}
\rhead{\includegraphics[scale=0.5]{logo_univ_rouen.png}}
\setlength{\headsep}{1cm}
\begin{document}

\maketitle
\newpage
\tableofcontents{}
\newpage
\section{Présentation}
  \paragraph{}
  Ce projet a pour but de comparer différentes bases de données NoSQL.
  Nous allons prendre les bases données CouchBase et MongoDB.
  
  
\section{La base de donnée CouchBase}
\subsection{Les modalités d’installation}
  \paragraph{} TODO
\subsection{Les méthodes d’insertion de données}
  \paragraph{} TODO
\subsection{Le langage de recherche}
  \paragraph{} TODO
\subsection{L’indexation interne}
  \paragraph{} TODO
\subsection{Le support de la concurrence d’accès}
  \paragraph{} TODO
\subsection{L’architecture du système}
  \paragraph{} TODO
\subsection{Les techniques de distribution}
  \paragraph{} TODO
  
  
  
  
\section{La base de donnée MongoDB}
\subsection{Les modalités d’installation}
  \paragraph{} TODO
\subsection{Les méthodes d’insertion de données}
  \paragraph{} TODO
\subsection{Le langage de recherche}
  \paragraph{} TODO
\subsection{L’indexation interne}
  \paragraph{} TODO
\subsection{Le support de la concurrence d’accès}
  \paragraph{} TODO
\subsection{L’architecture du système}
  \paragraph{} TODO
\subsection{Les techniques de distribution}
  \paragraph{} TODO


\section{Conclusion}
  \paragraph{} TODO
  
\end{document}